% ════════════════════════════════════════════
% 模型的改进与推广（8.模型改进推广.tex）写作规则 —— 模板内嵌，AI 先读本文件再写
% ════════════════════════════════════════════
\Needspace{10\baselineskip}
\section{模型的改进与推广}

\subsection{模型的改进}

后续改进可围绕当前尚未解释的差异、成本设定和数据覆盖展开。

在配比模型中，可对占比变量作适当变换，并加入领域交互项，再用 Group Lasso 等方法控制新增项数量，以考察 13 个领域平均 $R^2=0.629$ 之外约 $37\%$ 的变异。对于 nih\_exporter 等四个缺少验证损失的领域，可补充文本统计特征作为代理信息，减少固定参考占比的限制。冲突处理也可尝试有序 probit 因子模型，用共同的潜在质量解释指标分歧，进一步估计可信区间。

在质量模型中，可检验 $C(1-Q)^{\gamma_1}D^{-\gamma_2}$ 或其他质量与数据量交互形式，判断质量缺口是否会随数据规模变化。资源配置还可加入重复训练轮次、稀疏化与量化下的训练和推理费用，以及分阶段课程学习。针对低预算结果对成本形式较敏感的问题，更直接的办法是记录真实清洗流程的费用，拟合经验成本曲线后再与三类题设函数比较。

数据补充应优先服务于这些检验。可收集带质量标注的公开语料，扩展 B6/B7 目前覆盖的质量范围，并检查发布时间、许可、去重比例等信息能否辅助解释质量标签。C4 与 C1 可结合组织名、参数量和发布时间尝试匹配，经核实后再用于训练算力与能力的联合分析。若能获得 v1、v2 评测的共同校准样本，还可把 2019--2023 的历史记录换到同一尺度，减少当前可比窗口只有 10 个月的限制。

\subsection{模型的推广}

本文的方法可用于其他需要在有限成本下比较多项投入的场景，但应用前应重新检查数据、损失函数和成本关系。

在大模型预训练、微调和蒸馏中，可以沿用先拟合损失、再核算成本、最后优化配置的步骤。式\eqref{eq:p3_balance}和 $\tau(Q)$ 提供了比较语料处理与规模投入的方法；用于新任务时仍需重新估计对应参数，才能判断达到目标表现需要多少算力。

在语料治理中，22 个指标的方向表及冲突率相对独立参照的比较，可帮助开展文本筛选、清洗和配比调整。领域系数相似度 $s_{ij}$ 则可辅助寻找作用接近或不同的语料来源。具体能否等量替换、混合后是否有额外收益，仍应通过配比实验确认。

在科研资源规划中，可借鉴长周期 shift-share 分解与近期同一评测窗口对照的做法，分别观察扩参数与其他变化对应的增量。它能提供投入分析的线索，但非规模项不能全部归为算法收益。区分稳定结果、阶段边界和受假设影响较大的结果，也有助于多来源数据分析中的决策使用。

对于芯片制程、新能源和生物医药等可能逐渐趋于饱和的技术指标，也可比较 logistic 曲线与结构情景的预测差异，前提是先验证该指标确有相应增长特征。使用半合成或估算资料时，本文对评分方向和不可约损失下限的检查也可作为参考，帮助识别不能直接混合分析的数据。
